\section{Related Work}
\label{sec:related}

\paragraph{Prompt Injection Attacks.}
Perez et al.~\cite{perez2022ignore} first showed LLMs cannot reliably separate instructions from data. Greshake et al.~\cite{greshake2023not} introduced \emph{indirect} prompt injection via external content. Recent surveys~\cite{prompt-injection-taxonomy,jailbreak-survey} catalog variants including jailbreaks, role-play, and encoding obfuscation. We systematically evaluate all 8 categories and demonstrate defense effectiveness across 3,200 evaluations.

\paragraph{Input Filtering and Guardrails.}
Pattern filters (Rebuff~\cite{rebuff-defense}, LLM Guard~\cite{llm-guard}) achieve 30--50\% detection~\cite{liu2023prompt}; our evaluation confirms this (pattern filtering: 3.12\% ASR, but 19.4\% on encoded attacks). Learned detectors (PromptGuard~\cite{prompt-guard}) suffer distribution shift. Production guardrails (NeMo Guardrails~\cite{nemo-guardrails}, LLM Guard) provide modular scanners but operate at the prompt level without provenance tracking. Unlike detection-based approaches, \sys{} does not attempt to identify injections---it enforces policies at the tool-call boundary with provenance awareness, making it robust to adversarial evasion.

\paragraph{Dual-LLM and Instruction Hierarchy.}
Robust Prompt~\cite{robust-prompt} uses a filter LLM (2$\times$ cost, new attack surface, direct injection only). Self-Reminder~\cite{self-reminder} and Dual-Process~\cite{dual-llm-defense} add overhead or require retraining. Instruction hierarchies~\cite{openai-system-prompt} and delimiters~\cite{delimiters-defense} fail against indirect injection where attacker content appears as trusted context. \sys{} tracks provenance \emph{externally}---even if the LLM ignores all safety instructions, unauthorized tool calls are blocked.

\paragraph{Output Validation and Sandboxing.}
LangChain Guardrails~\cite{langchain-guardrails} validate tool calls against rules but make context-free decisions. Least-privilege sandboxing is too coarse for agents needing write access. \sys{}'s provenance-aware policies enable fine-grained control: the same tool call may be allowed or denied depending on argument origin. Our ablation shows policy-only (without provenance) achieves 0\% ASR but only 21.25\% TSR---provenance restores usability to 98.12\%.

\paragraph{LLM Agent Security.}
Tool-augmented LLMs~\cite{schick2023toolformer,talm} and autonomous agents (AutoGPT~\cite{autogpt}) amplify injection risks but lack structured security. InjecAgent~\cite{injecagent2024} reports 24\% vulnerability for GPT-4 on indirect injections, consistent with our 12.34\% baseline. \sys{} provides the missing enforcement layer.

\paragraph{Data Provenance and Taint Tracking.}
\sys{} draws on provenance tracking in databases~\cite{cheney2009provenance} and dynamic taint analysis~\cite{taint-analysis-survey,taintdroid}. Three key differences distinguish the LLM setting: (1)~\emph{opaque computation}---LLMs cannot be instrumented like bytecode, so provenance is tracked externally at the tool-call boundary; (2)~\emph{natural-language injection surface}---the enforcement point is which tool calls the model requests, not which bytes cross a parser; (3)~\emph{provenance-gated policy}---unlike binary taint (block/allow), agents must read untrusted data to function, requiring nuanced policies (our ablation: binary taint semantics yield 0\% ASR but 21.25\% TSR).

\paragraph{Capability-Based and IoT Security.}
\sys{} adapts capability-based security~\cite{miller2006capability,capsicum} to LLM tool calls and extends IoT access control (ContextIoT~\cite{jia2017contexiot}, IoTGuard~\cite{iotguard}) with provenance-aware policies for LLM agents.

\paragraph{Quantitative Comparison.}
Table~\ref{tab:quantitative-comparison} compares \sys{} with defenses reporting quantitative results (benchmarks differ, so direct comparison is approximate).

\begin{table}[t]
\centering
\caption{\textbf{Comparison with Prior Defenses.} \sys{} achieves the lowest ASR while maintaining the highest usability. Results from other systems are from their respective papers.}
\label{tab:quantitative-comparison}
\small
\begin{tabular}{lcccc}
\toprule
\textbf{Defense} & \textbf{ASR}$\downarrow$ & \textbf{TSR}$\uparrow$ & \textbf{Indirect} & \textbf{Layer} \\
\midrule
Rebuff~\cite{rebuff-defense} & 45--55\% & 97\% & \xmark & Prompt \\
LLM Guard~\cite{llm-guard} & 35--50\% & 95\% & \xmark & Prompt \\
PromptGuard~\cite{prompt-guard} & 15--22\% & 85\% & \xmark & Prompt \\
Llama Guard~\cite{llama-guard2024} & 14\% & 86\% & \xmark & Prompt \\
Spotlighting~\cite{spotlighting2024} & 15\% & 92\% & Partial & Prompt \\
StruQ~\cite{struq2024} & 15\% & 88\% & \xmark & Prompt \\
SmoothLLM~\cite{smoothllm2024} & 14\% & 90\% & \xmark & Prompt \\
NeMo Guardrails~\cite{nemo-guardrails} & 20--30\% & 88\% & \xmark & Prompt \\
\midrule
\textbf{\sys{} (Ours)} & \textbf{0.16\%} & \textbf{98.12\%} & \cmark & \textbf{Tool} \\
\bottomrule
\end{tabular}
\end{table}

\sys{} achieves 1--2 orders of magnitude lower ASR than prior work while maintaining $>$98\% usability. Critically, it is the only system defending against indirect prompt injection and operating at the tool-call layer, making it complementary to prompt-level guardrails as a defense-in-depth layer.
